% Generated by docs/generate_user_guide.py. Do not edit directly.
% Source: utilities/descriptions.yaml: approaches

\subsubsection{Approach 1: Direct Correlation}

\textbf{Fitting uncertainty}

Parameter estimates: fitted parameters are reported as estimate \(\pm\) 95\% confidence interval. The plus-minus value is calculated from the fitted model's covariance matrix, parameter standard errors, and residual degrees of freedom. It summarizes how precisely each parameter was estimated from the available data.

Fitting plot bands: shaded 95\% prediction bands are generated by Monte Carlo sampling of the fitted parameter covariance matrix. The sampled model variance is combined with the residual variance from the fitted data, then converted to a two-sided interval using a t multiplier when residual degrees of freedom are available.

These bands describe approximate prediction uncertainty around the fitted relationship, not regulatory acceptance limits. They depend on the fitted model form, local covariance approximation, residual behavior, and preprocessing choices.

\textbf{Prediction uncertainty}

Prediction plot error bars: when a prediction dataset is supplied, predicted in vivo times are plotted with horizontal error bars. The bars are calculated from the same 95\% prediction-band calculation used for the fitted correlation model; the upper and lower bar lengths are the distances from the band mean to the corresponding upper and lower interval bounds at each input prediction time.

The error bars therefore include propagated fitted-parameter uncertainty and residual/model error from the Approach 1 correlation fit.

\subsubsection{Approach 2: Independent Model Fitting}

\textbf{Fitting uncertainty}

Parameter estimates: each in vitro and in vivo model parameter is reported as estimate \(\pm\) 95\% confidence interval. The plus-minus value is calculated from the corresponding fitted model's covariance matrix, parameter standard errors, and residual degrees of freedom. It summarizes how precisely each parameter was estimated from the available data.

Time constants: tau values, when available, are calculated from the fitted model and parameter covariance. Reported tau values are shown as tau \(\pm\) 95\% confidence interval. The interval accounts for uncertainty in the fitted parameters used to calculate the time constant.

Fitting plot bands: the in vitro and in vivo shaded 95\% prediction bands are generated separately. For each fitted model, the app Monte Carlo samples the fitted parameter covariance matrix, combines sampled model variance with residual variance, and converts the resulting standard uncertainty to a two-sided interval using a t multiplier when residual degrees of freedom are available.

\textbf{Prediction uncertainty}

Response-value-ratio prediction error bars: uncertainty is propagated from both fitted models. The app samples each model's parameter covariance matrix, calculates the in vivo/in vitro response-value ratio for each sample, combines variance in that ratio with residual-based ratio uncertainty, and converts that uncertainty to a two-sided 95\% interval. The plotted vertical error bars are scaled from the resulting ratio interval at each prediction time.

Time-constant-ratio prediction error bars: when both tau values are valid, uncertainty from both tau values is carried through the ratio calculation. The resulting uncertainty in the scaled prediction time is converted to a symmetric 95\% interval, where possible, and plotted as a horizontal error bar.

Response-value ratio uncertainty can become unstable when the denominator model prediction or denominator time constant is near zero. Predictions in those regions should be interpreted cautiously.

\subsubsection{Approach 3: Direct Mapping}

\textbf{Fitting uncertainty}

Parameter uncertainty: Approach 3 does not fit a parametric model, so there is no fitted-parameter covariance matrix, no parameter confidence intervals, and no model-based 95\% prediction band.

Fitting plots: the plots show interpolation-based ratios/mappings calculated from the preprocessed datasets. They should be interpreted as deterministic mappings from the supplied data and preprocessing choices, not as statistical uncertainty bands.

\textbf{Prediction uncertainty}

Prediction plot error bars: Approach 3 does not generate formal propagated error bars for prediction outputs because no parametric covariance model is fitted.

Uncertainty is therefore limited to the quality of the input data, preprocessing, interpolation, and the assumption that the observed mapping remains valid for the prediction dataset. Extrapolated points outside the interpolation range are excluded or returned as unavailable.
